\documentclass[10pt,a4paper]{amsart}
\usepackage[margin=19mm]{geometry}
\usepackage{amsmath,amssymb,mathtools}
\usepackage[scr=rsfs,cal=euler]{mathalfa}
\usepackage{xfrac}
\usepackage[unicode,hidelinks,bookmarks=false]{hyperref}
\newcommand{\ie}{i.e.\ }
\newcommand{\cf}{cf.\ }
\newcommand{\resp}{respectively\ }
\newcommand{\Subsection}{\subsection}
\providecommand{\nobreakdash}{\nobreak\hbox{-}}
\setlength{\emergencystretch}{2em}
\multlinegap0pt
\usepackage{mathtools}
\usepackage{amssymb}
\usepackage{enumerate}
\newtheorem{theorem}{Theorem}
\newtheorem{lemma}{Lemma}
\title{Removable singularities for nonlocal minimal graphs}
\author{Minhyun Kim}
\date{Source: arXiv:2501.14299v2 (CC BY 4.0)}
\begin{document}
\maketitle

\noindent\emph{Source and licence.} Removable singularities for nonlocal minimal graphs. Version arXiv:2501.14299v2 (CC BY 4.0), distributed by arXiv under the Creative Commons Attribution 4.0 International licence, http://creativecommons.org/licenses/by/4.0/, as recorded in the arXiv licence metadata for this exact version and shown on the abstract page https://arxiv.org/abs/2501.14299v2. Full licence notice: \texttt{licenses/arxiv-2501.14299-CC-BY-4.0.txt}.\par\smallskip

\noindent\textit{Editorial scope.} Counted unit: the article's lemma lem-step2 (source0.tex lines 514--516). The article's complete original proof of it is delivered with it (source0.tex lines 519--590, one \texttt{proof} environment of 72 lines). No mathematics is written, repaired or re-derived: the statement and the proof are the article's own lines, in the article's own notation, and no retained line is substituted or re-flowed.  Retained uncounted as the source-local context the proof needs: lines 110--112; lines 228--242; lines 279--281; lines 298--312; lines 325--332; lines 334--341.  Nothing was omitted from any retained span, so the package carries no omission record.  No retained line contains a citation, so the wrapper carries no bibliography and no bibliography entry.  No retained line contains a figure environment, an image-inclusion command or drawing code, so no image is delivered and the package declares no asset dependency.  Editorial formatting only, in addition: an AMS \texttt{amsart} preamble that restates the article's own declarations and macros used by the retained lines, namely the environments \texttt{theorem}, \texttt{lemma} on separate counters, together with the standard packages those lines need.  The retained environments print in this document as Theorem 1, Lemma 1 in order of appearance, whereas the article prints the same items with the numbering of its own full document.  The complete native source of the article as distributed in the arXiv e-print for this exact version is bundled unchanged as \texttt{sources/172089/source0.tex}.
\begin{equation}\label{eq-main}
- \int_{\mathbb{R}^n} \mathscr{A}\left( x, y, \frac{u(x)-u(y)}{|x-y|} \right) \frac{\mathrm{d}y}{|x-y|^{n+s}} = \mathscr{B}(x, u),
\end{equation}

\begin{theorem}\label{thm-FSI}
Let $G \subset \mathbb{R}^n$ be a bounded open set with a Lipschitz boundary. Then, there exists a constant $C=C(n, s, G)>0$ such that
\begin{equation*}
\|u\|_{L^{1^{\ast}_s}(G)} \leq C \|u\|_{W^{s, 1}(G)}
\end{equation*}
for any $u \in W^{s, 1}(G)$, where
\begin{equation}\label{eq-1-ast}
1^{\ast}_s=\frac{n}{n-s}.
\end{equation}
In particular, there exists a constant $C=C(n, s)>0$ such that
\begin{equation*}
\|u\|_{L^{1^{\ast}_s}(B_r(x_0))} \leq C \left( [u]_{W^{s, 1}(B_r(x_0))} + r^{-s} \|u\|_{L^1(B_r(x_0))} \right)
\end{equation*}
for any $u \in W^{s, 1}(B_r(x_0))$.
\end{theorem}

\begin{lemma}\label{lem-cap-zero}
Let $E \subset \mathbb{R}^n$ be of $(s, 1)$-capacity zero. Then, there exists a sequence of functions $\bar{\eta}_j \in W^{s, 1}(\mathbb{R}^n)$ such that $\bar{\eta}_j=0$ in a neighborhood of $E$, $0 \leq \bar{\eta}_j \leq 1$ in $\mathbb{R}^n$, $\lim_{j\to \infty} [\bar{\eta}_j]_{W^{s, 1}(\mathbb{R}^n)} = 0$, and $\bar{\eta}_j \to 1$ a.e.\ in $\mathbb{R}^n$ as $j \to \infty$.
\end{lemma}

\begin{equation}\label{eq-truncation}
\bar{u}=
\begin{cases}
k &\text{if } u \leq k, \\
u &\text{if } k < u < l, \\
l &\text{if } u \geq l,
\end{cases}
\quad\text{and}\quad
\underbar{$u$}=
\begin{cases}
l &\text{if } u \leq -l, \\
-u &\text{if } -l < u < -k, \\
k &\text{if } u \geq -k.
\end{cases}
\end{equation}

\begin{equation}\label{eq-Caccio1-sub}
\begin{split}
&\int_{G} \int_{G} \frac{|\bar{v}^\beta(x) - \bar{v}^\beta(y)|}{|x-y|^{n+s}} (\eta\bar{\eta})(y) \,\mathrm{d}y \,\mathrm{d}x \\
&\leq C \int_{G} \int_{G} \bar{v}^\beta(x) \frac{|(\eta\bar{\eta})(x)-(\eta\bar{\eta})(y)|}{|x-y|^{n+s}} \,\mathrm{d}y \,\mathrm{d}x \\
&\quad + C \left( \beta (\mathrm{diam}\, G)^{1-s} + \sup_{x \in \mathrm{supp}\,\eta} \int_{\mathbb{R}^n \setminus G} \frac{\mathrm{d}y}{|x-y|^{n+s}} \right) \int_G \bar{v}^\beta \,\mathrm{d}x \\
&\quad + C \|f\|_{L^{n/s}(G \cap \{u>k\})} \|\bar{v}^\beta\eta\bar{\eta}\|_{L^{1^{\ast}_s}(G)}
\end{split}
\end{equation}

\begin{equation}\label{eq-Caccio1-super}
\begin{split}
&\int_{G} \int_{G} \frac{|\underbar{$v$}^\beta(x) - \underbar{$v$}^\beta(y)|}{|x-y|^{n+s}} (\eta\bar{\eta})(y) \,\mathrm{d}y \,\mathrm{d}x \\
&\leq C \int_{G} \int_{G} \underbar{$v$}^\beta(x) \frac{|(\eta\bar{\eta})(x)-(\eta\bar{\eta})(y)|}{|x-y|^{n+s}} \,\mathrm{d}y \,\mathrm{d}x \\
&\quad + C \left( \beta (\mathrm{diam}\, G)^{1-s} + \sup_{x \in \mathrm{supp}\,\eta} \int_{\mathbb{R}^n \setminus G} \frac{\mathrm{d}y}{|x-y|^{n+s}} \right) \int_G \underbar{$v$}^\beta \,\mathrm{d}x \\
&\quad + C \|f\|_{L^{n/s}(G \cap \{u<-k\})} \|\underbar{$v$}^\beta\eta\bar{\eta}\|_{L^{1^{\ast}_s}(G)},
\end{split}
\end{equation}

\begin{lemma}\label{lem-step2}
Let $\Omega \subset \mathbb{R}^n$ be open and let $K \subset \Omega$ be a compact set of $(s, 1)$-capacity zero. If $u$ is a solution of \eqref{eq-main} in $\Omega \setminus K$, then $u \in L^{1^{\ast}_s}_{\mathrm{loc}}(\Omega)$.
\end{lemma}

\begin{proof}
Let $\Omega' \Subset \Omega$ be an open set such that $K \subset \Omega'$. It is enough to show that $u \in L^{1^{\ast}_s}(\Omega')$. We fix a bounded open set $G$ with Lipschitz boundary such that $\Omega' \Subset G \Subset \Omega$. Let $\eta \in C^\infty_c(G)$ be a function such that $\eta =1$ on $\Omega'$, $0 \leq \eta \leq 1$, and $|\nabla \eta| \leq C$ for some $C>0$. By Lemma~\ref{lem-cap-zero}, there exist functions $\bar{\eta}_j \in W^{s, 1}(\mathbb{R}^n)$ such that $\bar{\eta}_j=0$ in a neighborhood of $K$, $0 \leq \bar{\eta}_j \leq 1$ in $\mathbb{R}^n$, $\lim_{j\to \infty} [\bar{\eta}_j]_{W^{s,1}(\mathbb{R}^n)} = 0$, and $\bar{\eta}_j \to 1$ a.e.\ in $\mathbb{R}^n$ as $j \to \infty$. Then, the estimate \eqref{eq-Caccio1-sub} with $\beta=1$ and $\bar{\eta}=\bar{\eta}_j$ shows that
\begin{align*}
I
&\coloneqq \int_{G} \int_{G} \frac{|\bar{u}(x) - \bar{u}(y)|}{|x-y|^{n+s}} (\eta\bar{\eta}_j)(y) \,\mathrm{d}y \,\mathrm{d}x \\
&\leq C_1 \int_{G} \int_{G} (\bar{u}(x)-k+\lambda) \frac{|(\eta\bar{\eta}_j)(x)-(\eta\bar{\eta}_j)(y)|}{|x-y|^{n+s}} \,\mathrm{d}y \,\mathrm{d}x \\
&\quad + C_1 \left( (\mathrm{diam}\, G)^{1-s} + \sup_{x \in \mathrm{supp}\,\eta} \int_{\mathbb{R}^n \setminus G} \frac{\mathrm{d}y}{|x-y|^{n+s}} \right) \int_{G} (\bar{u}-k+\lambda) \,\mathrm{d}x \\
&\quad + C_1 \|f\|_{L^{n/s}(G \cap \{u>k\})} \|(\bar{u}-k+\lambda)\eta\bar{\eta}_j\|_{L^{1^{\ast}_s}(G)},
\end{align*}
where $\bar{u}$ is defined as in \eqref{eq-truncation} and $C_1=C_1(n, s, \Lambda, \lambda)>0$.

By the fractional Sobolev inequality in Theorem~\ref{thm-FSI}, we have that
\begin{align*}
\|(\bar{u}-k+\lambda)\eta\bar{\eta}_j\|_{L^{1^{\ast}_s}(G)}
&\leq C_2 \int_{G} \int_{G} (\bar{u}(x)-k+\lambda) \frac{|(\eta\bar{\eta}_j)(x)-(\eta\bar{\eta}_j)(y)|}{|x-y|^{n+s}} \,\mathrm{d}y \,\mathrm{d}x \\
&\quad + C_2 I + C_2 \|(\bar{u}-k+\lambda)\eta\bar{\eta}_j\|_{L^1(G)}
\end{align*}
for some $C_2=C_2(n, s, G)>0$. If we take $k$ sufficiently large so that
\begin{equation}\label{eq-k}
C_1C_2 \|f\|_{L^{n/s}(G \cap \{u>k\})} \leq 1/2,
\end{equation}
then
\begin{equation*}
\frac{1}{2} I \leq C \int_{G} \int_{G} (\bar{u}(x)-k+\lambda) \frac{|(\eta\bar{\eta}_j)(x)-(\eta\bar{\eta}_j)(y)|}{|x-y|^{n+s}} \,\mathrm{d}y \,\mathrm{d}x + C \int_G (\bar{u}-k+\lambda) \,\mathrm{d}x
\end{equation*}
for some $C>0$, where we computed
\begin{equation*}
\sup_{x \in \mathrm{supp}\,\eta} \int_{\mathbb{R}^n \setminus G} \frac{\mathrm{d}y}{|x-y|^{n+s}} \leq \sup_{x \in \mathrm{supp}\,\eta} \int_{\mathbb{R}^n \setminus B_d(x)} \frac{\mathrm{d}y}{|x-y|^{n+s}} \leq \frac{|\mathbb{S}^{n-1}|}{s d^s}
\end{equation*}
with $d=\mathrm{dist}(\mathrm{supp}\,\eta, \partial G)>0$. Another application of the fractional Sobolev inequality and the triangle inequality
\begin{equation*}
|(\eta\bar{\eta}_j)(x)-(\eta\bar{\eta}_j)(y)| \leq |\bar{\eta}_j(x)-\bar{\eta}_j(y)| + |\eta(x)-\eta(y)|
\end{equation*}
yield that
\begin{align*}
&\|(\bar{u}-k+\lambda)\eta\bar{\eta}_j\|_{L^{1^{\ast}_s}(G)} \\
&\leq C \int_{G} \int_{G} (\bar{u}(x)-k+\lambda) \frac{|(\eta\bar{\eta}_j)(x)-(\eta\bar{\eta}_j)(y)|}{|x-y|^{n+s}} \,\mathrm{d}y \,\mathrm{d}x + C \int_{G} (\bar{u}-k+\lambda) \,\mathrm{d}x \\
&\leq C (l-k+\lambda) [\bar{\eta}_j]_{W^{s, 1}(G)} + C \int_{G} (\bar{u}-k+\lambda) \,\mathrm{d}x.
\end{align*}
Since $\lim_{j\to \infty} [\bar{\eta}_j]_{W^{s,1}(\mathbb{R}^n)} = 0$, by letting $j \to \infty$, we have that
\begin{align*}
\|\bar{u}-k\|_{L^{1^{\ast}_s}(\Omega')}
&\leq \|\bar{u}-k+\lambda\|_{L^{1^{\ast}_s}(\Omega')} \\
&\leq \|(\bar{u}-k+\lambda)\eta\|_{L^{1^{\ast}_s}(G)} \\
&\leq C_3 \|\bar{u}-k+\lambda\|_{L^1(G)} \\
&\leq C_3 \|\bar{u}-k\|_{L^1(\Omega')} + C_3 \|\bar{u}-k\|_{L^1(G\setminus \Omega')} + C_3 \lambda |G|
\end{align*}
for some $C_3>0$. Observing that $\bar{u}-k=0$ in $\{u \leq k\}$ and using H\"older's inequality, we obtain that
\begin{equation*}
\|\bar{u}-k\|_{L^1(\Omega')} = \|\bar{u}-k\|_{L^1(\Omega' \cap \{u > k\})} \leq |\Omega' \cap \{u>k\}|^{s/n} \|\bar{u}-k\|_{L^{1^{\ast}_s}(\Omega')}.
\end{equation*}
By taking $k$ sufficiently large so that \eqref{eq-k} and $C_3 |\Omega' \cap \{u>k\}|^{s/n} \leq 1/2$ hold, we have that
\begin{equation*}
\|\bar{u}-k\|_{L^{1^{\ast}_s}(\Omega')} \leq C \|\bar{u}-k\|_{L^1(G\setminus \Omega')} + C
\end{equation*}
for some $C>0$. Letting $l$ tend to infinity shows that
\begin{equation*}
\|u-k\|_{L^{1^{\ast}_s}(\Omega' \cap \{u>k\})} \leq C \|u-k\|_{L^1((G\setminus \Omega')\cap \{u>k\})} + C,
\end{equation*}
and consequently,
\begin{align}\label{eq-step2}
\begin{split}
\|u_+\|_{L^{1^{\ast}_s}(\Omega')}
&\leq C \|u\|_{L^{1^{\ast}_s}(\Omega' \cap \{u>k\})} + C \|u\|_{L^{1^{\ast}_s}(\Omega' \cap \{0<u\leq k\})} \\
&\leq C \|u-k\|_{L^{1^{\ast}_s}(\Omega' \cap \{u>k\})} + C \|k\|_{L^{1^{\ast}_s}(\Omega' \cap \{u>k\})} + C \|k\|_{L^{1^{\ast}_s}(\Omega')} \\
&\leq C \|u-k\|_{L^1((G\setminus \Omega')\cap \{u>k\})} + C + 2Ck|\Omega'|^{1-s/n}.
\end{split}
\end{align}
Since $u \in L^1_{\mathrm{loc}}(\Omega \setminus K)$ and $G \setminus \Omega' \Subset \Omega \setminus K$, the right-hand side of \eqref{eq-step2} is finite.

Similarly, one can show that $u_- \in L^{1^{\ast}_s}(\Omega')$ by using the estimate \eqref{eq-Caccio1-super} instead of \eqref{eq-Caccio1-sub}. This concludes that $u \in L^{1^{\ast}_s}(\Omega')$.
\end{proof}

\end{document}
